\section{Pair profiles and the regular-set estimates}
\label{p14:comp:pair-profiles}

This section supplies the estimates used in Section~\ref{p1:arith:section}.
The regular-set transfer is used for the explicit channel formulas and for
the positional comparison. The scalar reconstruction with exact marginals
does not discard these exceptional windows.

\subsection{Regular-set transfer}
\label{p14:comp:regular-transfer}

\begin{proof}[Proof of Proposition~\ref{p1:arith:core}]
The one-window estimates proved in Lemma~\ref{p1:arith:defects} give
$m_B(x)=O_C(B/\log B)$ uniformly, and only $N^{o(1)}$ windows contain two
$B$-defective integers.  A nonzero relative relation is an even nonempty
selection.  Private prime coordinates therefore force it to use at least two
defects.  Its dimension is at most $m_B(x)$.  Consequently
\[
 \sum_{x\in I_N}(2^{L-\dim S_x}-1)=N^{o(1)}.
\]
Fourier inversion proves \eqref{p1:arith:marginalsum}.  The sum of the
probabilities of active discarded starts is at most $pN^{o(1)}$; the coupling
using the original prime signs proves the total-variation assertion.  On $G_N$
the marginal is exactly $p$, giving both assertions about the mean.

For the last statement, let $D=\{Z_N\ne Z_N^\circ\}$.  Since
$0\leq Z_N^\circ\leq Z_N$, H\"older's inequality gives, for any fixed $s>1$,
\[
 \E\big[Z_N^j\ind_D\big]
 \leq(\E Z_N^{js})^{1/s}\P(D)^{1-1/s}.
\]
Choose a sufficiently large integer $s$, and then use the fixed-moment bound.
This last step is used only after that bound has been established; a small
coupling error alone would not justify a variance transfer.
\end{proof}

\subsection{Sparse hosts and weighted residual mass}
\label{p14:comp:sparse-hosts}

\begin{proof}[The elementary bounds preceding Proposition~\ref{p1:arith:hosts}]
A relative relation that is internal to one window contributes at most
$N^{1+o(1)}$ hosts, by Proposition~\ref{p1:arith:core}.  Otherwise it selects
at least two vertices in each window.  Put $T=N^{3/5}$.  If all selected
kernels are at most $T$, each window contains two small kernels.  The estimate
\begin{equation}\label{p1:arith:twokernelused}
 \#\{x\in I_N:\#\{n\in V_x:K_B(n)\leq N^\beta\}\geq2\}
   \ll_{C,\beta,\varepsilon}N^{\max(\beta,7/12)+\varepsilon}
\end{equation}
proved in Lemma~\ref{p1:arith:twokernel} gives $N^{6/5+\varepsilon}$
choices.  If one selected kernel exceeds $T$, Lemma~\ref{p1:arith:anchor}
gives
$N^{3/2+o(1)}T^{-1/2+\tau}$, for every fixed $\tau>0$.
Choose $\tau$ sufficiently small.

For the weighted assertion, first restrict to $c(x,y)<\eta B$.
Equation~\eqref{p1:arith:graphbound} bounds the weight by
$2^{\eta B+o(B)}$; choose $\eta$ after the desired error exponent.
The complementary sector has total weight $O_{C,\eta,\varepsilon}(N^{1+\varepsilon})$
by Lemma~\ref{p1:arith:dense}.  The latter estimate uses the component
equations themselves, rather than just their support: many small components
put the pair on a multiquadratic surface, and its low-degree curves are either
aligned or absent from the dyadic region.
This proves both bounds with exponent $6/5+\varepsilon$, which is all
that the pair-profile proof below uses. The exponent $1+\varepsilon$
in the revised statement is proved later in
Section~\ref{p14:comp:confinement}; no late confinement estimate is
being inferred from this elementary calculation.
\end{proof}

\subsection{The complete pair profile}
\label{p14:comp:complete-profile}

\begin{proof}[Proof of Theorem~\ref{p1:arith:pairs}]
For the channel $2:3$, either orientation permits
$N/2\leq t<2N/3$, up to $O(B)$ endpoint changes.  For $3:4$ the corresponding
interval is $N/3\leq t<N/2$.  Both intervals have length $N/6+O(B)$.
The two orientations therefore give $N/3$ times the finite sums above.
Translations near a boundary contribute $O(B^3 2^{L/b})$.

For either fixed channel choose $H$ a sufficiently large fixed multiple of
$B$ and normalize the values to $bV_x\cup aV_y$.  Its diameter is $O(B)$.
If its distinct values are all nondefective at level $H$, each has a private
prime and the only relations before normalization are the even selections
of exact units.  Lemma~\ref{p1:arith:defects} shows that the exceptional
channel pairs number $N^{1/2+o(1)}$.  Lemma~\ref{p1:arith:alignedrank} bounds
their nullity by $L/b+o(L)$.  Replacing the rational kernel by the full
arithmetic kernel consequently costs at most
$N^{1/2+1/b+o(1)}$, for either the positive or the signed sum.

There are $N B^{O(1)}$ aligned hosts altogether.  The slope of height two
forces both starts within $O(B)$ of the dyadic endpoints; its mass is
$N^{1/2+o(1)}$.  Height one has no exact unit in separated windows.
The only reduced slopes of heights three and four that occur in the interior
are $2:3$ and $3:4$.  Every remaining channel has height at least five and
nullity at most $L/5+o(L)$, again by Lemma~\ref{p1:arith:alignedrank}; its
total contribution is $N^{6/5+o(1)}$. The elementary version just
proved controls the nonaligned sector at the same sufficient scale.
The inequality following
\eqref{p1:arith:affine} permits exactly the same error estimates for the signed
sum.

At distances $1,\ldots,L-1$, two run starts are incompatible.  At distance
$L$, their edge systems form a tree on an interval of diameter $2L$;
Lemma~\ref{p1:arith:defects}, at threshold $2L+1$, bounds its nullity by
$o(L)$.  Hence the entire contact band changes the second factorial moment
by $O(N^{-1+o(1)})$.  Multiplying \eqref{p1:arith:signedprofile} by $p^2$,
and using $\mu_N-\lambda=O(N^{-1+o(1)})$, proves
\eqref{p1:arith:cumulantexpansion}.  Finally
$2^{(L-r)/b}=2^{-r/b}(N/\lambda)^{1/b}$ in
\eqref{p1:arith:S23}--\eqref{p1:arith:S34} gives the two equivalents.
\end{proof}

\subsection{The good-edge mass}
\label{p14:comp:good-edge-mass}

\begin{proof}[Proof of Corollary~\ref{p1:arith:edgefamily}]
Each start has $O(B)$ possible partners in a fixed channel.  The rational
kernel has dimension one less than the number of exact units when that
number is positive, giving the rank bound.  All pairs excluded from the good
family have precisely the error mass estimated in the proof of
Theorem~\ref{p1:arith:pairs}, together with the contact band.
For good compatible pairs the centered contribution is $q_e-p^2$.
Since $p^2|\mathcal E_N|=O(B/N)$, replacing it by $q_e$ has the stated cost.
The asymptotic follows from \eqref{p1:arith:leadingphase}.
\end{proof}

\subsection{A further window after a good edge}
\label{p14:comp:pair-palm}

The following channel-specific estimate is used in the first reconstruction
and in the positional comparison. Its good-edge notation is that of
Section~\ref{p1:arith:section}.

\begin{proposition}[Positive Palm cost of a good edge]\label{p1:arith:palm}
For $a:b\in\{2:3,3:4\}$ let $\mathcal E_{a,b}$ be its good edges and put
$\beta_{a,b}=\sum_{e\in\mathcal E_{a,b}}q_e$.  Then
\begin{align}
 \frac1{\beta_{2,3}}\sum_{\{x,y\}\in\mathcal E_{2,3}}q_{xy}
    \sum_{z\in G_N\setminus\{x,y\}}
       (\P(J_z=1\mid J_x=J_y=1)-p)_+
   &\ll_{C,\varepsilon}N^{-2/3+\varepsilon},\label{p1:arith:palm23}\\
 \frac1{\beta_{3,4}}\sum_{\{x,y\}\in\mathcal E_{3,4}}q_{xy}
    \sum_{z\in G_N\setminus\{x,y\}}
       (\P(J_z=1\mid J_x=J_y=1)-p)_+
   &\ll_{C,\varepsilon}N^{-1/2+\varepsilon}.\label{p1:arith:palm34}
\end{align}
Consequently the unnormalised positive two-block masses are
$O_{C,\varepsilon}(N^{-4/3+\varepsilon})$ and
$O_{C,\varepsilon}(N^{-5/4+\varepsilon})$, respectively.
\end{proposition}

\begin{proof}
For a compatible good source $e=\{x,y\}$ and $z\notin e$,
\[
 q_e\bigl(\P(J_z=1\mid J_x=J_y=1)-p\bigr)_+
       =\bigl(q(e\cup\{z\})-q_ep\bigr)_+.
\]
A positive summand has separated starts and positive crossing above
the source. It is at most $W(e\cup\{z\})$, so part (2) of
Proposition~\ref{p1:arith:fixed-homogeneous-seeds} bounds the total.
The source phase gives $\beta_{a,b}\asymp_{C,a,b}LN^{-1+1/b}>0$.
Divide by this positive quantity and reduce the internal exponent
margin. This gives both displayed normalized bounds and the stated
unnormalised bounds.
\end{proof}
